% =====================================================================
% laporan/template-laporan.tex
% Kerangka laporan akhir Final Project Konsep Basis Data.
% Setiap bagian diberi komentar "% ISI:" yang menyebut bab panduan sumber
% bahannya, supaya penulis tidak mulai dari halaman kosong.
% Bangun dengan: bash tools/build.sh
% =====================================================================
\documentclass[12pt,a4paper]{article}

\usepackage[T1]{fontenc}
\usepackage[utf8]{inputenc}
\usepackage{lmodern}
\usepackage[indonesian]{babel}
\usepackage[a4paper,margin=2.5cm,top=3cm]{geometry}
\usepackage{microtype}
\usepackage{booktabs}
\usepackage{tabularx}
\usepackage{longtable}
\usepackage{array}
\usepackage{enumitem}
\usepackage{xcolor}
\usepackage{listings}
\usepackage{titlesec}
\usepackage{graphicx}
\usepackage{tikz}
\usetikzlibrary{shapes.geometric}
\usepackage[most]{tcolorbox}
\usepackage[hidelinks]{hyperref}

\definecolor{sqlkata}{RGB}{0,90,160}
\definecolor{sqlkomentar}{RGB}{90,120,90}
\definecolor{kotakgaris}{RGB}{70,110,150}

\lstdefinestyle{sql}{
  language=SQL,
  basicstyle=\ttfamily\footnotesize,
  keywordstyle=\color{sqlkata}\bfseries,
  commentstyle=\color{sqlkomentar}\itshape,
  numbers=left,numberstyle=\tiny\color{gray},numbersep=6pt,
  frame=single,rulecolor=\color{kotakgaris},
  breaklines=true,showstringspaces=false,tabsize=2
}
\lstset{style=sql}

% Kotak untuk menempelkan tangkapan layar hasil eksekusi.
\newtcolorbox{tangkapanlayar}[1]{
  colback=gray!6,colframe=kotakgaris,boxrule=0.6pt,arc=2pt,
  fonttitle=\bfseries\small,title={Tangkapan layar: #1},
  height=6cm,valign=center,halign=center
}

\titleformat{\section}{\large\bfseries}{BAB \Roman{section}}{0.6em}{\centering}
\titleformat{\subsection}{\normalsize\bfseries}{\thesubsection}{0.5em}{}
\renewcommand{\thesection}{\Roman{section}}
\renewcommand{\thesubsection}{\arabic{section}.\arabic{subsection}}
\setlist{nosep,leftmargin=1.5em}

\begin{document}

% ---------------------------------------------------------------------
% Sampul laporan
% ISI: data kelompok. Ganti bagian bergaris bawah dengan nama dan NIM.
% ---------------------------------------------------------------------
\begin{titlepage}
  \centering
  \vspace*{1.5cm}
  {\Large\textbf{LAPORAN FINAL PROJECT}}\\[0.5em]
  {\large KONSEP BASIS DATA}\\[2em]

  {\LARGE\textbf{Sistem Informasi Klinik Pratama}}\\[1em]
  {\large Perancangan dan Implementasi Basis Data Relasional}\\[3em]

  \begin{tabular}{ll}
    Disusun oleh: & \\[0.6em]
    Anggota 1 & \underline{\hspace{5cm}} \\
              & NIM \underline{\hspace{3cm}} \\[0.4em]
    Anggota 2 & \underline{\hspace{5cm}} \\
              & NIM \underline{\hspace{3cm}} \\[0.4em]
    Anggota 3 & \underline{\hspace{5cm}} \\
              & NIM \underline{\hspace{3cm}} \\[0.4em]
    Anggota 4 & \underline{\hspace{5cm}} \\
              & NIM \underline{\hspace{3cm}} \\[0.4em]
    Anggota 5 & \underline{\hspace{5cm}} \\
              & NIM \underline{\hspace{3cm}} \\[0.4em]
    Anggota 6 & \underline{\hspace{5cm}} \\
              & NIM \underline{\hspace{3cm}} \\
  \end{tabular}\\[2.5em]

  {\large Program Studi S1 Informatika}\\
  {\large Fakultas Ilmu Komputer}\\
  {\large Universitas Pembangunan Nasional Veteran Jakarta}\\
  {\large 2026}
\end{titlepage}

\tableofcontents
\newpage

% =====================================================================
\section{Deskripsi Basis Data}
% =====================================================================

\subsection{Latar belakang dan studi kasus}

Klinik Pratama adalah klinik rawat jalan berskala kecil dengan lima poli: Poli
Umum, Poli Anak, Poli Gigi, Poli Penyakit Dalam, dan Poli Mata. Setiap poli
punya dokter sendiri, dan seorang dokter bertugas di satu poli. Dokter praktik
pada sesi tertentu, misalnya tanggal 5 Agustus 2026 pukul 08.00 sampai 12.00 di
ruang R-01 dengan kuota sepuluh pasien.

Pelayanan berjalan dalam lima langkah. Pasien mendaftar di meja depan pada satu
sesi jadwal, lalu menerima nomor antrean. Dokter memeriksa dan mencatat keluhan,
diagnosis, tekanan darah, suhu, berat badan, serta tinggi badan. Setelah itu
dokter memberi tindakan dan menulis resep. Di meja pembayaran, petugas
menghitung tagihan dari biaya tindakan ditambah harga obat, lalu mencatat
pembayaran. Setiap langkah menghasilkan data yang harus tersimpan pada tempat
yang tepat, karena langkah berikutnya membaca hasil langkah sebelumnya.

Basis data relasional dipilih karena tiga hal pada alur itu menuntut penjagaan
mesin. Nomor antrean tidak boleh ganda pada satu jadwal, jumlah pasien tidak
boleh melewati kuota, dan satu kunjungan hanya boleh punya satu baris
pembayaran. Ketiganya bisa dinyatakan sebagai kunci unik, kunci tamu, atau
trigger, sehingga pelanggaran ditolak database tanpa bergantung pada disiplin
petugas. Pertimbangan kedua, laporan pendapatan per poli dan pemakaian obat
membutuhkan penggabungan beberapa tabel yang sudah dinormalisasi. Pertimbangan
ketiga, harga obat dan tarif tindakan berubah dari waktu ke waktu, dan rujukan
ke satu baris master menjaga nilai lama tetap konsisten.

\subsection{Ruang lingkup}

Rancangan mencakup pendaftaran pasien, penjadwalan sesi praktik, kunjungan,
rekam medis, tindakan, resep beserta stok obat, tagihan, pembayaran, pelaporan,
hak akses, indeks, dan pengujian. Empat hal sengaja diletakkan di luar cakupan
supaya batas pekerjaan jelas.

\begin{table}[ht]
\centering
\caption{Hal di luar cakupan beserta alasannya}
\small
\begin{tabularx}{\textwidth}{@{}p{4.2cm}X@{}}
\toprule
\textbf{Di luar cakupan} & \textbf{Alasan} \\
\midrule
Riwayat cicilan pembayaran & Pembayaran dianggap lunas sekali bayar, sehingga status dan tanggal sudah memadai \\
Pengadaan obat dari pemasok & Stok hanya berkurang lewat resep; penambahan dilakukan manual dengan SQL \\
Jadwal berulang mingguan & Setiap sesi dicatat satu baris pada tanggal tertentu, nama hari diturunkan dengan \texttt{DAYNAME()} \\
Autentikasi aplikasi & Hak akses diatur di tingkat database memakai \texttt{GRANT}, tanpa aplikasi \\
\bottomrule
\end{tabularx}
\end{table}

Dua hal lain juga tidak diambil. Jenis kelamin pasien tidak dipakai satu pun
kebutuhan informasi, dan aplikasi antarmuka tidak dibuat karena penilaian
terfokus pada perancangan serta SQL.

\subsection{Aturan bisnis}

Aturan bisnis ditulis sebagai pernyataan yang bisa diperiksa mesin. Kalimat
seperti ``pelayanan harus ramah'' tidak dipakai karena tidak punya penjaga
teknis. Setiap aturan di bawah ini disertai penjaga dan pengujian yang
membuktikannya.

\begin{longtable}{@{}p{0.7cm}p{6.1cm}p{4.4cm}p{3.2cm}@{}}
\caption{Delapan aturan bisnis, penjaga, dan bukti pengujian}
\label{tab:aturan}\\
\toprule
\textbf{No} & \textbf{Aturan} & \textbf{Penjaga} & \textbf{Bukti} \\
\midrule
\endfirsthead
\toprule
\textbf{No} & \textbf{Aturan} & \textbf{Penjaga} & \textbf{Bukti} \\
\midrule
\endhead
1 & Kuota jadwal tidak boleh terlampaui & trigger \texttt{trg\_cek\_kuota\_jadwal} & Uji 4 \\
2 & Satu dokter tidak boleh punya dua jadwal beririsan pada tanggal yang sama & trigger \texttt{trg\_cek\_bentrok\_jadwal} & Uji 5 \\
3 & Satu pasien hanya boleh terdaftar sekali pada satu sesi jadwal & \texttt{uq\_kunjungan\_pasien} & \texttt{verify.sh} bagian E \\
4 & Resep hanya boleh dibuat setelah rekam medis ada & kunci tamu \texttt{fk\_detail\_resep\_rekam} & skema \\
5 & Satu kunjungan punya paling banyak satu baris pembayaran & \texttt{uq\_pembayaran\_kunjungan} & Uji 8 \\
6 & Total tagihan adalah jumlah biaya tindakan ditambah harga obat & \texttt{sp\_bayar\_kunjungan} menghitung sendiri & Uji 7 \\
7 & Stok obat berkurang otomatis saat resep disimpan & \texttt{trg\_cek\_stok\_obat} dan \texttt{trg\_kurangi\_stok} & Uji 6 \\
8 & Status kunjungan hanya salah satu dari empat nilai & \texttt{ENUM} pada \texttt{status\_kunjungan} & skema \\
\bottomrule
\end{longtable}

Aturan 1 menjaga jumlah pasien per sesi, sedangkan aturan 7 menjaga jumlah obat
per resep. Keduanya bekerja pada tabel berbeda dan diuji di skenario berbeda.
Penolakan atas pelanggaran aturan 1, 2, dan 7 muncul sebagai galat 1644 dari
trigger, sedangkan pelanggaran aturan 3, 5, dan 8 ditolak lewat kunci dan
\texttt{ENUM}.

\subsection{Peran dan tanggung jawab anggota}

Pembagian peran mengikuti aturan satu berkas satu pemilik. Kelompok berjumlah
enam anggota; NIM diisi setelah data kelompok final, sedangkan nama dan uraian
pekerjaan sudah tetap.

\begin{table}[ht]
\centering
\caption{Pembagian peran anggota kelompok}
\small
\begin{tabularx}{\textwidth}{@{}p{0.8cm}p{3.4cm}X@{}}
\toprule
\textbf{No} & \textbf{Nama dan NIM} & \textbf{Tanggung jawab} \\
\midrule
1 & Agit & Studi kasus, kamus data, ERD, delapan kebutuhan informasi, Bab I dan Bab III laporan, perbandingan NoSQL, \texttt{sql/00\_kesepakatan.sql}, \texttt{sql/04\_prosedur.sql} \\
2 & Dapa & \texttt{sql/02\_data.sql}, dokumentasi normalisasi 1NF sampai BCNF, tangkapan layar query dasar \\
3 & Farras & \texttt{sql/01\_struktur.sql}, tabel transformasi, skema relasional, contoh aljabar relasional \\
4 & Kenzie & \texttt{sql/03\_view\_trigger.sql}, \texttt{sql/06\_indeks.sql}, penjelasan view, trigger, dan indeks \\
5 & Kafka & \texttt{sql/05\_keamanan.sql}, \texttt{sql/08\_kueri.sql}, penjelasan hak akses dan sepuluh query bukti \\
6 & Rassya & \texttt{sql/07\_pengujian.sql}, skrip \texttt{tools/}, berkas \texttt{bukti/}, demo isolasi, daftar periksa, build PDF \\
\bottomrule
\end{tabularx}
\end{table}

Review silang berjalan dua arah. Farras memeriksa hasil Kenzie dan sebaliknya,
Kafka memeriksa hasil Agit dan Dapa, sedangkan Agit memeriksa daftar periksa
buatan Rassya. Perubahan nama tabel atau kolom hanya lewat Farras, dan
perubahan angka data contoh hanya lewat Dapa.

\subsection{Kamus data}

Kamus data adalah daftar resmi nama entitas dan atributnya. Semua berkas SQL
memakai nama pada Tabel \ref{tab:kamusdata}, sehingga tidak ada singkatan baru
yang muncul di tengah pengerjaan. Tipe data dan nama constraint tidak diulang di
sini karena sumbernya \texttt{sql/01\_struktur.sql}.

\begin{longtable}{@{}p{3.1cm}p{5.4cm}p{7.4cm}@{}}
\caption{Kamus data sebelas entitas}
\label{tab:kamusdata}\\
\toprule
\textbf{Entitas} & \textbf{Atribut} & \textbf{Keterangan} \\
\midrule
\endfirsthead
\toprule
\textbf{Entitas} & \textbf{Atribut} & \textbf{Keterangan} \\
\midrule
\endhead
poli & id\_poli, nama\_poli, lantai & Tempat layanan. Nama poli unik. \\
dokter & id\_dokter, id\_poli, nama\_dokter, no\_sip, email, no\_telp, tipe\_dokter & Tenaga medis. Nomor SIP dan email unik. \\
jadwal\_praktik & id\_jadwal, id\_dokter, tanggal, jam\_mulai, jam\_selesai, ruangan, kuota & Satu baris untuk satu sesi pada satu tanggal. \\
pasien & id\_pasien, nama\_pasien, nik, email, no\_telp, tanggal\_lahir, alamat & NIK 16 digit dan unik. \\
kunjungan & id\_kunjungan, id\_pasien, id\_jadwal, tanggal\_kunjungan, no\_antrean, status\_kunjungan & Pendaftaran pasien pada satu sesi jadwal. \\
rekam\_medis & id\_rekam, id\_kunjungan, keluhan, diagnosis, tekanan\_darah, suhu, berat\_badan, tinggi\_badan, imt, tanggal\_periksa & Hasil pemeriksaan. Satu kunjungan paling banyak satu rekam. \\
tindakan & id\_tindakan, nama\_tindakan, tarif & Master tindakan dan tarifnya. \\
detail\_tindakan & id\_detail\_tindakan, id\_kunjungan, id\_tindakan, jumlah & Tindakan yang dikerjakan pada satu kunjungan. \\
obat & id\_obat, nama\_obat, satuan, harga, stok & Master obat dan sisa stok. \\
detail\_resep & id\_detail\_resep, id\_rekam, id\_obat, jumlah, aturan\_pakai & Obat yang diresepkan pada satu rekam medis. \\
pembayaran & id\_pembayaran, id\_kunjungan, tanggal\_bayar, metode, total\_tagihan, status\_bayar & Tagihan satu kunjungan. Satu kunjungan paling banyak satu baris. \\
\bottomrule
\end{longtable}

\subsection{Delapan kebutuhan informasi}

Kebutuhan informasi dirumuskan sebagai pertanyaan supaya bisa langsung diuji
dengan query. Kolom terakhir menyebut query yang menjawabnya pada
\texttt{bukti/kueri-panduan.sql}.

\begin{table}[ht]
\centering
\caption{Kebutuhan informasi dan query yang menjawabnya}
\small
\begin{tabularx}{\textwidth}{@{}p{0.7cm}Xp{2.2cm}@{}}
\toprule
\textbf{No} & \textbf{Pertanyaan} & \textbf{Query} \\
\midrule
1 & Pasien mana yang lahir sebelum 2005, diurutkan dari yang paling tua? & Q1 \\
2 & Dokter mana bertugas di poli apa, dan poli itu di lantai berapa? & Q2 \\
3 & Bagaimana riwayat kunjungan lengkap: pasien, dokter, poli, status? & Q3 \\
4 & Berapa sisa kuota setiap jadwal praktik saat ini? & Q4 \\
5 & Kunjungan mana yang tagihannya belum lunas? & Q5 \\
6 & Pasien mana yang pernah dirujuk ke Poli Anak? & Q6 \\
7 & Poli mana yang pendapatannya di atas rata-rata? & Q7 \\
8 & Obat apa yang paling banyak keluar, dan dokter mana yang menangani kunjungan terbanyak? & Q8, Q10 \\
\bottomrule
\end{tabularx}
\end{table}

Dua pertanyaan pertama berhenti pada satu atau dua tabel. Sisanya menuntut
penggabungan tiga tabel atau lebih, dan pertanyaan nomor 8 menuntut dua bentuk
agregasi berbeda dalam satu kebutuhan.

% =====================================================================
\section{Perancangan dan Implementasi}
% =====================================================================

\subsection{Diagram relasi entitas}

Gambar \ref{gam:erd-laporan} menampilkan sebelas entitas dan relasi di
antaranya. Atribut tidak digambar supaya diagram tetap terbaca pada ukuran
halaman; daftar atribut lengkap ada pada kamus data di Tabel
\ref{tab:kamusdata}.

\begin{figure}[ht]
\centering
\resizebox{0.98\textwidth}{!}{\input{../panduan/gambar/erd}}
\caption{Diagram relasi entitas Sistem Informasi Klinik Pratama}
\label{gam:erd-laporan}
\end{figure}

Delapan relasi pada diagram punya tiga bentuk kardinalitas. Relasi 1:N menjadi
kunci tamu di sisi banyak, relasi 1:0..1 menjadi kunci tamu unik di sisi anak,
dan relasi M:N menjadi tabel penghubung tersendiri. Tabel
\ref{tab:kardinalitas-laporan} menuliskan kedelapan relasi beserta bentuk tabel
hasilnya.

\begin{table}[ht]
\centering
\caption{Kardinalitas relasi dan bentuk tabel hasil}
\label{tab:kardinalitas-laporan}
\small
\begin{tabularx}{\textwidth}{@{}p{4.5cm}p{1.6cm}X@{}}
\toprule
\textbf{Relasi} & \textbf{Kardinalitas} & \textbf{Bentuk tabel hasil} \\
\midrule
poli memiliki dokter & 1 : N & kunci tamu \texttt{id\_poli} di \texttt{dokter} \\
dokter punya jadwal praktik & 1 : N & kunci tamu \texttt{id\_dokter} di \texttt{jadwal\_praktik} \\
pasien mendaftar kunjungan & 1 : N & kunci tamu \texttt{id\_pasien} di \texttt{kunjungan} \\
jadwal diisi kunjungan & 1 : N & kunci tamu \texttt{id\_jadwal} di \texttt{kunjungan} \\
kunjungan menghasilkan rekam medis & 1 : 0..1 & \texttt{id\_kunjungan} unik di \texttt{rekam\_medis} \\
kunjungan menerima tindakan & M : N & tabel penghubung \texttt{detail\_tindakan} \\
kunjungan ditagih pembayaran & 1 : 0..1 & \texttt{id\_kunjungan} unik di \texttt{pembayaran} \\
rekam medis diresepkan obat & M : N & tabel penghubung \texttt{detail\_resep} \\
\bottomrule
\end{tabularx}
\end{table}

Delapan relasi itu menghasilkan sepuluh kunci tamu. Jumlahnya lebih besar dari
jumlah relasi karena dua relasi 1:N menuju tabel yang sama, yaitu
\texttt{kunjungan}, sehingga tabel itu memuat dua kunci tamu sekaligus.

Satu catatan model extended ER berlaku pada entitas \texttt{dokter}. Model ER
dasar tidak bisa menyatakan bahwa sebagian dokter adalah spesialis dan sebagian
dokter umum. Model extended ER menggambarkannya sebagai spesialisasi, dengan
dokter sebagai entitas induk dan dua subentitas di bawahnya. Implementasi di
sini memilih satu tabel dengan kolom pembeda \texttt{tipe\_dokter} bertipe
\texttt{ENUM('Umum','Spesialis')}. Atribut yang berbeda antar subentitas masih
sedikit, sehingga memecah tabel menjadi tiga hanya menambah penggabungan tanpa
manfaat nyata. Jika nanti dokter spesialis memerlukan belasan atribut sendiri,
tabel terpisah baru masuk akal.

\subsection{Transformasi ke skema relasional}
% ISI: ambil dari Bab 7 panduan.
% Tuliskan sebelas relasi hasil transformasi beserta kunci primer dan
% kunci tamunya. Sebutkan aturan transformasi yang dipakai untuk relasi
% 1:N, M:N, dan 1:0..1.

\subsection{Normalisasi}
% ISI: ambil dari Bab 8 panduan.
% Jelaskan tahap unnormalized, 1NF, 2NF, 3NF, dan BCNF. Sertakan tabel
% dependensi fungsional dan bukti bahwa setiap determinan adalah
% candidate key.

\subsection{Source code}
% ISI: salin isi sql/01_struktur.sql sampai sql/07_pengujian.sql.
% Kode tidak perlu seluruhnya; tampilkan bagian penting per berkas,
% misalnya definisi tabel kunjungan, satu view, satu trigger, dan satu
% prosedur.

\subsubsection{Struktur tabel}
\begin{lstlisting}
-- ISI: tempel CREATE TABLE untuk tiga tabel inti, misalnya
-- poli, kunjungan, dan rekam_medis, dari sql/01_struktur.sql.
\end{lstlisting}

\subsubsection{View dan trigger}
\begin{lstlisting}
-- ISI: tempel definisi v_kunjungan_lengkap dan trg_cek_kuota_jadwal
-- dari sql/03_view_trigger.sql.
\end{lstlisting}

\subsubsection{Stored procedure}
\begin{lstlisting}
-- ISI: tempel definisi sp_bayar_kunjungan dari sql/04_prosedur.sql.
\end{lstlisting}

\subsubsection{Hak akses dan indeks}
\begin{lstlisting}
-- ISI: tempel CREATE USER, GRANT, dan CREATE INDEX dari
-- sql/05_keamanan.sql dan sql/06_indeks.sql.
\end{lstlisting}

\subsection{Demo program}
% ISI: ambil dari Bab 10, 11, 12, dan 13 panduan.
% Setiap tangkapan layar diberi keterangan satu baris yang menyebut
% perintah yang dijalankan dan hasil yang terlihat.

\subsubsection{Data yang sudah dimuat}
\begin{tangkapanlayar}{isi tabel kunjungan setelah reset-db.sh}
  Tempel tangkapan layar di sini.
\end{tangkapanlayar}

\subsubsection{Hasil query lanjutan}
\begin{tangkapanlayar}{hasil query Q3 dan Q4}
  Tempel tangkapan layar di sini.
\end{tangkapanlayar}

\subsubsection{Penolakan aturan bisnis}
\begin{tangkapanlayar}{pesan galat kuota jadwal sudah penuh}
  Tempel tangkapan layar di sini.
\end{tangkapanlayar}

\subsubsection{Transaksi dan tingkat isolasi}
\begin{tangkapanlayar}{keluaran tools/demo-isolasi.sh}
  Tempel tangkapan layar di sini.
\end{tangkapanlayar}

\subsubsection{Hak akses dan indeks}
\begin{tangkapanlayar}{hasil SHOW GRANTS dan EXPLAIN}
  Tempel tangkapan layar di sini.
\end{tangkapanlayar}

\subsubsection{Pengujian otomatis}
\begin{tangkapanlayar}{keluaran tools/verify.sh yang seluruhnya lulus}
  Tempel tangkapan layar di sini.
\end{tangkapanlayar}

% =====================================================================
\section{Kesimpulan dan Saran}
% =====================================================================

\subsection{Kesimpulan}

Rancangan ini menjawab delapan kebutuhan informasi yang ditetapkan pada Bab I.
Sebagian pertanyaan berhenti pada satu atau dua tabel, sedangkan pertanyaan
tentang riwayat kunjungan dan pendapatan poli menuntut penggabungan lebih dari
empat tabel. Riwayat kunjungan lengkap, misalnya, baru terjawab setelah
\texttt{kunjungan}, \texttt{pasien}, \texttt{jadwal\_praktik}, \texttt{dokter},
dan \texttt{poli} digabungkan, dan pertanyaan pendapatan per poli menambah
\texttt{pembayaran} ke dalam penggabungan itu. Kesepuluh query menghasilkan
angka yang sama dengan \texttt{bukti/10-kueri.txt}, termasuk 14 baris riwayat
kunjungan, 3 tagihan belum lunas, dan Amlodipine 5mg sebanyak 60 tablet sebagai
obat terbanyak.

Sebelas tabel dan sepuluh kunci tamu sudah cukup mewakili alur klinik dari
pendaftaran sampai pembayaran. Rancangan tidak menambah tabel keempat belas
untuk keperluan yang belum ada, dan tidak menyimpan nilai turunan yang bisa
dihitung. Indeks massa tubuh dihitung database, nama hari diambil dari tanggal,
dan total tagihan dibentuk saat pembayaran dijalankan.

Aturan bisnis ditegakkan mesin, bukan diingat petugas. \texttt{verify.sh}
menjalankan 34 pemeriksaan dan berakhir dengan 34 lulus, 0 gagal. Pemeriksaan
bagian E mencatat penolakan nyata: nomor antrean ganda ditolak dengan galat
1062, pasien yang tidak terdaftar ditolak 1452, suhu 50 derajat ditolak 3819,
pendaftaran melewati kuota dan jadwal bentrok ditolak 1644, begitu pula resep
yang melebihi stok.

Transaksi pembayaran memenuhi sifat seluruhnya atau tidak sama sekali.
Pemanggilan pembayaran kedua ditolak dengan pesan \texttt{Kunjungan sudah
dibayar}, dan percobaan insert yang gagal tidak menyisakan baris. Demo isolasi
menunjukkan perbedaan dua tingkat isolasi pada kasus yang sama: pada READ
COMMITTED pembacaan kedua melihat stok 195, sedangkan pada REPEATABLE READ
pembacaan kedua tetap 95 setelah sesi lain mengubah nilainya, sesuai aturan
tingkat isolasi tersebut.

Hak akses dan indeks juga terbukti bekerja. Akun \texttt{petugas\_klinik}
ditolak saat mencoba menulis ke tabel \texttt{obat} dengan galat 1142. Setelah
\texttt{idx\_pasien\_nama} dibuat, rencana eksekusi berubah dari pemindaian
penuh 10 baris menjadi pencarian rentang 2 baris; pada
\texttt{idx\_kunjungan\_tanggal} jumlah baris yang diperiksa turun dari 14
menjadi 1.

\subsection{Kendala}

Pemeriksaan \texttt{CHECK} baru ditegakkan MySQL mulai versi 8.0.16. Kelompok
memutuskan memakai versi itu atau lebih baru pada seluruh mesin, karena
menjalankan skrip pada versi lama membuat aturan rentang suhu, berat, dan tinggi
tidak berfungsi tanpa pesan galat apa pun.

Data contoh terikat pada 34 pemeriksaan sekaligus pada tabel hasil di panduan
dan berkas \texttt{bukti/}. Satu angka yang berubah menuntut
\texttt{tools/reset-db.sh}, \texttt{tools/verify.sh}, dan seluruh keluaran bukti
dijalankan ulang. Karena itu penyuntingan \texttt{sql/02\_data.sql} dibatasi
pada Anggota 2.

Efek indeks pada tabel kecil terbatas. Tabel \texttt{pasien} hanya berisi 10
baris dan \texttt{kunjungan} 14 baris, sehingga pemindaian penuh pun masih
cepat. Perbedaan yang bisa dibuktikan ada pada rencana eksekusi dan jumlah baris
yang diperiksa, bukan pada waktu tanggap.

\subsection{Saran pengembangan}

Pengadaan obat dari pemasok bisa ditambahkan sebagai tabel batch yang memuat
nomor batch, tanggal kedaluwarsa, dan jumlah masuk. Konsekuensinya, kolom
\texttt{stok} pada \texttt{obat} tidak lagi cukup sebagai satu angka, karena
sisa stok harus dijumlahkan dari batch yang belum kedaluwarsa.

Jadwal berulang mingguan bisa ditambahkan sebagai tabel pola jadwal ditambah
tabel pengecualian tanggal. Trigger pemeriksa bentrok kemudian harus memeriksa
sesi hasil perluasan pola, bukan hanya baris yang tersimpan langsung.

Riwayat cicilan bisa ditambahkan sebagai tabel pembayaran rinci yang merujuk
satu baris \texttt{pembayaran}. Dengan begitu \texttt{total\_tagihan} berubah
menjadi nilai turunan dari jumlah cicilan, dan status lunas dihitung dari sisa
tagihan.

Dua saran operasional menyusul: pencadangan otomatis harian dengan
\texttt{mysqldump}, dan laporan bulanan terjadwal yang membaca
\texttt{v\_pendapatan\_poli} sehingga petugas tidak perlu menjalankan query
manual.

\subsection{Perbandingan dengan pendekatan NoSQL}

Materi minggu 15 dinilai lewat pemetaan struktur, bukan pembangunan server
NoSQL. Satu kunjungan beserta rekam medis dan resepnya bisa ditulis sebagai satu
dokumen JSON seperti berikut.

\begin{lstlisting}[numbers=none]
{
  "id_kunjungan": 8,
  "tanggal_kunjungan": "2026-09-10",
  "no_antrean": 1,
  "status_kunjungan": "Selesai",
  "pasien": { "id_pasien": 8, "nama_pasien": "Hesti Purwanti" },
  "jadwal": {
    "id_jadwal": 5,
    "tanggal": "2026-09-10",
    "jam_mulai": "08:00:00",
    "jam_selesai": "11:00:00",
    "dokter": { "id_dokter": 5, "nama_dokter": "dr. Eka Lestari, Sp.M" },
    "poli": { "nama_poli": "Poli Mata", "lantai": "3" }
  },
  "rekam_medis": {
    "keluhan": "Mata merah dan berair",
    "diagnosis": "Konjungtivitis bakterial",
    "tekanan_darah": "112/72",
    "suhu": 37.1,
    "imt": 21.5,
    "resep": [
      { "id_obat": 7, "nama_obat": "Salep Mata Kloramfenikol",
        "jumlah": 1, "aturan_pakai": "Oles tipis 3x sehari" }
    ]
  },
  "pembayaran": { "metode": "Tunai", "total_tagihan": 140000.00,
                  "status_bayar": "Lunas" }
}
\end{lstlisting}

Data yang hampir selalu dibaca bersama sebaiknya dilekatkan. Rekam medis dan
resepnya selalu dibuka bersamaan saat dokter menelusuri riwayat, jadi keduanya
masuk ke dalam dokumen kunjungan dan tidak perlu penggabungan tambahan. Data
yang dipakai banyak dokumen sebaiknya dirujuk. Harga obat dan tarif tindakan
dipakai ribuan kunjungan, sehingga menyalin nilainya ke setiap dokumen membuat
satu perubahan harga harus diperbarui di semua tempat. Pada dokumen di atas,
obat tetap dirujuk lewat \texttt{id\_obat}, sedangkan nama obat hanya disalin
untuk keperluan tampilan.

Skema relasional tetap lebih tepat untuk klinik ini, dan alasannya bukan
kebiasaan memakai SQL. Alasan pertama, aturan bisnis di sini harus ditegakkan
mesin, dan kunci unik, kunci tamu, \texttt{CHECK}, serta transaksi sudah
melakukannya tanpa kode tambahan. Dokumen JSON tidak punya penjaga setara untuk
nomor antrean ganda atau pembayaran ganda; keduanya harus diperiksa di lapisan
aplikasi. Alasan kedua, sebagian besar kebutuhan informasi berupa agregasi
lintas poli dan lintas dokter, misalnya pendapatan poli di atas rata-rata atau
peringkat dokter. Agregasi semacam itu diselesaikan satu query pada tabel
ternormalisasi, sedangkan pada dokumen menuntut pemrosesan seluruh koleksi.
Alasan ketiga, harga obat dan tarif tindakan berubah dan nilai lama harus tetap
konsisten pada transaksi yang sudah terjadi. Rujukan ke satu baris master
menjamin hal itu, sementara penyalinan nilai ke dokumen membuka kemungkinan dua
dokumen memuat harga berbeda untuk obat yang sama.

% =====================================================================
% Daftar pustaka
% ISI: sesuaikan dengan sumber yang benar-benar dipakai. Cantumkan
% modul kuliah dan dokumentasi resmi MySQL.
% =====================================================================
\begin{thebibliography}{9}
  \bibitem{modul} Tim Pengampu Konsep Basis Data. \textit{Modul Kuliah Konsep
    Basis Data}. Program Studi S1 Informatika, UPN Veteran Jakarta, 2026.
  \bibitem{mysql} Oracle Corporation. \textit{MySQL 8.0 Reference Manual}.
    \url{https://dev.mysql.com/doc/refman/8.0/en/}, diakses 2026.
  \bibitem{elm} Elmasri, R. dan Navathe, S. B. \textit{Fundamentals of Database
    Systems}. Edisi ke-7. Pearson, 2016.
\end{thebibliography}

\end{document}
